%%%%%%%%%%%%%%%%%%%%%%%%%%%%%%%%%%%%%%%%%%%%%%%%%%%%%%%%%%%%%%%
%
% Welcome to Overleaf --- just edit your LaTeX on the left,
% and we'll compile it for you on the right. If you open the
% 'Share' menu, you can invite other users to edit at the same
% time. See www.overleaf.com/learn for more info. Enjoy!
%
%%%%%%%%%%%%%%%%%%%%%%%%%%%%%%%%%%%%%%%%%%%%%%%%%%%%%%%%%%%%%%%
\documentclass[10pt]{article}
\usepackage[a4paper, total={8in, 11.4in} ]{geometry}
\usepackage{graphicx}
\usepackage{amssymb}

\renewcommand{\baselinestretch}{1}

\title{}

\begin{document}
\fontsize{14}{10}\textbf{\underline{Definitions:}}
\par

\fontsize{12}{10}
\textbf{\textit{\underline{Week 1:}}}
\textbf{\underline{Definition 1.1:}}
Let A and B be two sets. We say that A is a subset of B, or that A is contained
in B if and only if each x $ \in $A is also an element of B. We write  A $ \subseteq $ B to say that “A is a
subset of B.” If there is an element of A not in B, we write A $ \nsubseteq $ B.
\par


\textbf{\underline{Definition 1.2:}}
The empty set, written $ \emptyset $, is the set that has no elements.
\par


\textbf{\underline{Definition 1.3:}}
If A and B are statements, their conjunction is the statement “A and B”.
A $ \wedge $ B is true exactly when both A is true and B is true
\par
\textbf{\underline{Definition 1.4:}}
The intersection S$\cap $T is the set consisting of the
elements S and T have in common. One can write S$ \cap $T=${x:(x \in S)\wedge(x \in T)}$
\par
\textbf{\underline{Definition 1.5:}}
The \textit{disjunction} of two statements A and B is the statement “A or B”. A $\wedge $ B is true exactly 
when either A is true or B is true.
\par
\textbf{\underline{Definition 1.6:}}
The union S $ \cap $ T is the set consisting of the elements
of either S or T. One can write S$\cap$T= ${x:(x\in S)\vee(x\in T)}$
\par
\textbf{\underline{Definition 1.7:}}
Let T$ \subseteq $S be a subset.The complement of T in S is
the subset S $ \setminus $ T =${x\in S:x \notin T} \subseteq S $
\par
\textbf{\underline{Definition 1.8:}}
The converse to the conditional A $ \Rightarrow $B is the statement B$ \Rightarrow $A.
\par
\textbf{\underline{Definition 1.9:}}
The contrapositive of a conditional A $ \Rightarrow $ B is the statement  $ \overline{\mbox{B}} $  $ \Rightarrow $ $\overline{\mbox{A}}$ 
\par
\textbf{\underline{Definition 1.10:}}
We say that two statements A and B are equivalent to mean that A is true
if and only if B is true.
\par
\textbf{\underline{Definition 1.11:}}
\par
\textbf{\underline{Definition2.1:}}
A field is a set F along with operations + and · such that the following hold for a,b,c $\in$ F:
\underline{Addition:} \underline{A(0)}: a+b $\in$ F, \underline{A(1)}:Commutativity, a+b=b+a, \underline{A(2)}: Associativity, a+(b+c)=(a+b)+c, \underline{A(3)}: There is an element 0 $\in$ F such that a+0=a
for all a $\in$ F, \underline{A(4)}: There is an element -a $\in$ F such that a + (-a)=0.
\underline{Multiplication:} \underline{M(0)}: a $\cdot$ b $\in$ F. \underline{M(1)} Commutativity a·b = b·a \underline{M(2)} Associativity a·(b·c)=(a·b)·c
\underline{M(4)} if a$\neq$0 there exists number 1/a where a·1/a =1. \underline{M(5)} Distributivity,a·(b+c)=(a·b)+(a·c)
\par
\textbf{\textit{\underline{Week 2:}}}
\textbf{\underline{Definition 3.1:}}
An ordered field is a field F, together with a relation $<$, so that the following
hold: \underline{(R1)}  (Positive/negative) If x $\in$ F, then exactly one of the following is true: 0 $<$ x, x =0,
or x$<$0. \underline{(R2)}(Negations reverse) If y $<$x, then -x$<$-y. \underline{(R3)}(Can add constants) If x$<$y, and c$\in$F, then x+c$<$y+c.
\underline{(R4)} (Positives multiply) If 0 $<$ x and 0$<$y then 0$<$xy. \underline{(R5)} (Transitive property) If x $<$ y and y$<$z then x$<$z.
\textbf{all but 2 is minimal}
\par
\textbf{\underline{Definition 3.2:}}
For x $\in$ R, we define the absolute value or modulus of x to be
$|x| =$ $ \{ \begin{array}{rcl}  
        x & \mbox{if} & x \geq  0 \\
        -x & \mbox{if} & x \leq  0 
        \end{array} $ 
\par
\textbf{\underline{Definition 4.1:}}
Given a sequence of numbers a$_m$,a$_{m+1}$  \dots  a$_n$ we write
 $\sum_{k=m}^{n} a_k = a_m +a_{m+1} + \dots a_n$
\par
\textbf{\underline{Definition 4.2:}}
\textbf{The Principle of Mathematical Induction}
Given a list of statements
P(k), P(k + 1),..., we may conclude that P(n) is true for every integer n $\geq$ k provided that
(1) we know that P(k) is true ("base case") (2) and we can prove that P(n) $\Rightarrow$ P(n + 1) for any integer n $\geq$ k ("induction step").
\par
\textbf{\textit{\underline{Week 3:}}}
\textbf{\underline{Definition 5.1:}}
Let A be a subset of R An upper bound for A is a real number M such that for all x $\in$ A, we have x $\leq$ M.
We say that A is bounded above if there exists an upper bound for A, and unbounded
above otherwise.
• A lower bound for A is a real number m such that for all x $\in$ A, we have m $\leq$ x. We
say that A is bounded below if there exists a lower bound for A, and unbounded below
otherwise.
• We say A is bounded if it is both bounded above and bounded below.
\par
\textbf{\underline{Definition 5.2:}}
Given a subset A $\subseteq$ R, a number L is a least upper bound (LUB) or \textit{supremum}
for A if
(a) L is an upper bound for A, that is, x $\leq$ L for all x $\in$ A; and
(b) L $\leq$ M for every upper bound M of A.
Similarly, we say $\ell$ is a greatest lower bound (GLB) or \textit{infimum} for A if
(a) $\ell$ is a lower bound for A, that is, x $\geq$ $\ell$ for all x $\in$ A; and
(b) $\ell$ $\geq$ m for every lower bound m of A.
\end{document}